\chapter{Stability and level-set geometry}\label{sec:phase-stability}
A positive linearized solution, the geometry of regular level sets, and a growth bound give different routes to rigidity. We develop \cite{ChAC}, Proposition 4.1 and Sections 5.1--5.4. Elliptic estimates and unique continuation are identified where used; the local computations do not replace those analytic inputs.

\begin{definition}[Local Allen-Cahn stability]\label{def:phase-local-stability}
A smooth solution of $\varepsilon^2\Delta_g u=W'(u)$ on an open Riemannian manifold is stable if
\[
 Q_\varepsilon(\phi)=\int\left(\varepsilon|\nabla\phi|^2+
 \varepsilon^{-1}W''(u)\phi^2\right)dV_g\geq0
 \quad\text{for every }\phi\in C_c^\infty.
\]
On a closed manifold all smooth tests are allowed. Compactly supported $H^1$ tests are also admissible by approximation, since the potential is bounded on their common compact support. This is the local form of Proposition \ref{prop:allen-cahn-variation}, as used in \cite{ChAC}, Section 5.1.
\end{definition}

\begin{proposition}[A positive linearized solution gives stability]\label{prop:phase-ground-state}
Let $u$ be a smooth Allen--Cahn solution. If $v>0$ is smooth and satisfies $\varepsilon^2\Delta_g v=W''(u)v$, then for every compactly supported smooth $\phi$,
\[
 Q_\varepsilon(\phi)=\varepsilon\int v^2
 \left|\nabla\left(\frac\phi v\right)\right|^2dV_g\geq0.
\]
\end{proposition}
\begin{proof}
Write $\phi=v\eta$ and expand its gradient. The terms other than $v^2|\nabla\eta|^2$ integrate to
$\int\langle\nabla v,\nabla(v\eta^2)\rangle+
\varepsilon^{-2}W''(u)v^2\eta^2=0$ by the equation for $v$. There is no boundary term because $\eta$ is compactly supported and $v$ stays positive there. Multiply by $\varepsilon$. This is the same ground-state mechanism as Proposition \ref{prop:positive-jacobi-stability}, applied to the Allen--Cahn operator; compare \cite{ChAC}, Lemma 5.6.
\end{proof}

\begin{corollary}[Monotone entire solutions are stable]\label{cor:phase-monotone-stability}
If a smooth Euclidean Allen--Cahn solution has $\partial_eu>0$ everywhere for one fixed unit vector $e$, it is stable for compactly supported variations. In particular every planar heteroclinic is stable.
\end{corollary}
\begin{proof}
Differentiate the equation in the constant direction $e$. The positive function $v=\partial_eu$ satisfies the linearized equation, so Proposition \ref{prop:phase-ground-state} applies. A planar heteroclinic has a strictly positive derivative in its transition direction. This proves \cite{ChAC}, Lemma 5.6 without imposing the additional limits needed for the minimizing theorem below.
\end{proof}

\begin{lemma}[The gradient defect and regular level geometry]\label{lem:phase-level-geometry}
Let $u$ be smooth, $w=|\nabla u|$, and $\nu=\nabla u/w$ on the regular set $O=\{w>0\}$. If $A$ is the second fundamental form of a regular level set with normal $\nu$, then on $O$,
\[
 |\nabla^2u|^2-|\nabla w|^2
 =w^2|\nabla\nu|^2
 =w^2|A|^2+|\nabla_{\{u=\mathrm{constant}\}}w|^2\geq0.
\]
Also $w$ is locally Lipschitz and $|\nabla w|\leq|\nabla^2u|$ almost everywhere on the whole manifold.
\end{lemma}
\begin{proof}
Differentiate $\nabla u=w\nu$. The derivatives of the unit vector are perpendicular to it, so the squared norm gives $|\nabla^2u|^2=|\nabla w|^2+w^2|\nabla\nu|^2$. In a frame tangent to the level set, the tangential derivatives of $\nu$ are its second fundamental form. Symmetry of the Hessian gives, for tangent $X$, $w\langle\nabla_\nu\nu,X\rangle=Xw$, so the remaining normal-direction contribution is $w^{-2}|\nabla_{\{u=\mathrm{constant}\}}w|^2$.

The norm of a smooth vector field is locally Lipschitz. The usual norm chain rule gives the inequality where $w>0$; on its zero set the weak derivative of $w$ vanishes almost everywhere. Thus the a.e. inequality holds globally. This proves the identities behind \cite{ChAC}, Exercise 5.5(b),(c),(f).
\end{proof}

\begin{theorem}[The geometric stability inequality]\label{thm:phase-geometric-stability}
For a stable smooth Allen--Cahn solution on a Riemannian manifold and every $\eta\in C_c^\infty$,
\[
 \int\left(|\nabla^2u|^2-|\nabla|\nabla u||^2+
 \operatorname{Ric}(\nabla u,\nabla u)\right)\eta^2
 \leq\int|\nabla u|^2|\nabla\eta|^2.
\]
\end{theorem}
\begin{proof}
Put $w=|\nabla u|$. Bochner and the equation give
\[
 \tfrac12\Delta_g w^2=|\nabla^2u|^2+
 \varepsilon^{-2}W''(u)w^2+\operatorname{Ric}(\nabla u,\nabla u).
\]
The compactly supported $H^1$ test $\eta w$ is admissible in stability. Expanding $|\nabla(\eta w)|^2$, use $w\nabla w=\tfrac12\nabla w^2$ to integrate the cross term by parts. After division by $\varepsilon$, the potential term cancels the same term in the Bochner formula. The result is exactly the asserted inequality. All products are locally integrable; the norm test can equivalently be approximated by $\eta\sqrt{w^2+\delta^2}$ and then $\delta\downarrow0$. This proves \cite{ChAC}, Proposition 4.1 and Exercise 4.2, with the test-function justification explicit.
\end{proof}

\begin{corollary}[Stable phases on a closed positive-Ricci manifold]\label{cor:phase-positive-ricci}
On a connected closed Riemannian manifold with positive Ricci curvature, the only smooth stable Allen--Cahn solutions are $u\equiv+1$ and $u\equiv-1$, for every $\varepsilon>0$.
\end{corollary}
\begin{proof}
Use $\eta=1$ in Theorem \ref{thm:phase-geometric-stability}. The Hessian defect is nonnegative by Lemma \ref{lem:phase-level-geometry}; positivity of Ricci therefore forces $\nabla u=0$. Connectedness makes $u$ constant, and the equation permits only $-1,0,1$. The value zero is unstable because $W''(0)=-1$ and the constant test gives $Q_\varepsilon(1)=-\varepsilon^{-1}\operatorname{Vol}(M)<0$. At the wells $W''(\pm1)=2$, so stability is immediate. This is the closed-manifold form of \cite{ChAC}, Exercise 4.3(b).
\end{proof}

\begin{proposition}[Quadratic gradient growth makes the level geometry flat]\label{prop:phase-quadratic-flatness}
Let $u:\mathbb R^n\to[-1,1]$, $n\geq2$, be a stable solution with $\varepsilon=1$. Suppose
\[
 \int_{B_R}|Du|^2\leq CR^2\quad(R\geq1).
\]
Then $|D^2u|^2-|D|Du||^2=0$ on $\{|Du|>0\}$. Consequently the unit gradient direction is constant on each connected regular component, and its level sets there have zero second fundamental form.
\end{proposition}
\begin{proof}
For $R=2^k$, use the logarithmic cutoff $\eta_R=1$ on $B_R$, $\eta_R=\log(R^2/|x|)/\log R$ on $B_{R^2}\setminus B_R$, and zero outside. Compactly supported Lipschitz tests follow by approximation. The geometric stability inequality bounds the nonnegative defect integral against $\eta_R^2$ by
\[
 \frac1{(\log R)^2}\int_{B_{R^2}\setminus B_R}|x|^{-2}|Du|^2
 \leq\frac{4Ck}{(k\log2)^2}\longrightarrow0.
\]
The estimate sums over $k$ dyadic annuli, using the gradient-growth assumption on the enclosing ball. Fatou gives zero defect almost everywhere; continuity on the regular set gives zero everywhere there. Lemma \ref{lem:phase-level-geometry} gives $D\nu=0$, and connectedness makes its value constant. This is the weighted logarithmic-cutoff argument in \cite{ChAC}, Section 5.4, and parallels Proposition \ref{prop:quadratic-growth-cutoffs}.
\end{proof}

\begin{lemma}[From zero level defect to a one-dimensional solution]\label{lem:phase-flat-to-onedimensional}
Let $u:\mathbb R^n\to[-1,1]$ be a nonconstant smooth solution with $\varepsilon=1$. If the defect in Lemma \ref{lem:phase-level-geometry} vanishes on its regular set, then $u(x)=f(e\cdot x)$ for a unit vector $e$ and a smooth function $f:\mathbb R\to[-1,1]$.
\end{lemma}
\begin{proof}
Choose a connected regular component. There the unit gradient is one constant vector $e$. For each constant $a\perp e$, the function $v=\partial_a u$ vanishes on this nonempty open component and satisfies $\Delta v=W''(u)v$ on all of $\mathbb R^n$. Unique continuation for this scalar linear elliptic equation gives $v=0$ everywhere. This unique-continuation theorem is an external analytic input, as in \cite{ChAC}, Exercise 5.5(d); it is not proved by the preceding energy estimate. Choosing a basis of $e^\perp$ shows that $u$ is constant on every affine hyperplane perpendicular to $e$, hence has the asserted form.
\end{proof}

\begin{lemma}[Bounded stable one-dimensional profiles]\label{lem:phase-stable-ode}
A bounded $C^2$ solution $f:\mathbb R\to[-1,1]$ of $f''=W'(f)$ is stable on $\mathbb R$ if and only if it is a constant well or a translated, possibly reversed heteroclinic. If a solution $u(x)=f(e\cdot x)$ is stable on $\mathbb R^n$, then $f$ is stable on $\mathbb R$.
\end{lemma}
\begin{proof}
For the second assertion rotate $e$ to the last coordinate and test with $\phi(t)\eta(y/R)$, where $\phi$ is compactly supported and $\eta\not\equiv0$ is a fixed compactly supported transverse cutoff. Divide the stability inequality by $\int\eta(y/R)^2$. The transverse gradient contribution tends to zero like $R^{-2}$, giving the one-dimensional stability inequality. The case $n=1$ needs no reduction.

The constant zero is unstable by testing with $\phi(t/L)$ and taking $L$ large: its negative potential term is of order $L$ and the gradient term of order $L^{-1}$. A nonconstant solution cannot attain $\pm1$ or attain $0$ with zero derivative, by ODE uniqueness. If it has a turning point with value $a\in(-1,1)\setminus\{0\}$, its conserved Hamiltonian gives a periodic orbit between $-|a|$ and $|a|$. Indeed $f'^2=2(W(f)-W(a))$, positive between the two turning values, and the separated travel-time integrals are finite because their endpoint roots are simple.

Let $(s,t)$ lie between successive extrema of this periodic orbit, and extend $h=f'$ on $[s,t]$ by zero. Then $h\in H^1_c(\mathbb R)$, $h(s)=h(t)=0$, and integration of its linearized equation gives $Q(h)=0$. Nonnegative quadratic forms obey $Q(h,\zeta)=0$ for all compact tests, by applying stability to $h+\tau\zeta$ for both signs of $\tau$. But integration by parts gives
$Q(h,\zeta)=h'(t)\zeta(t)-h'(s)\zeta(s)$, and $h'=f''=W'(f)\ne0$ at either extremum. Choosing a test supported near just one endpoint contradicts this identity. Thus a nonconstant stable profile has no turning point.

It is strictly monotone and bounded, so it has distinct limits at both ends. Uniform derivative bounds and convergence of translates show that these limits are equilibria in $\{-1,0,1\}$, with derivative tending to zero. Conservation of the Hamiltonian requires equal potential values at the two ends; the only distinct pair is $-1,+1$. Separation of variables now gives the heteroclinic. Conversely constant wells are stable, and heteroclinics are stable by their positive translation derivative and Proposition \ref{prop:phase-ground-state}. This resolves the ODE step in \cite{ChAC}, Exercise 5.5(e).
\end{proof}

\begin{theorem}[Stable rigidity under quadratic gradient growth]\label{thm:phase-stable-growth-rigidity}
Under the hypotheses of Proposition \ref{prop:phase-quadratic-flatness}, either $u\equiv\pm1$ or
\[
 u(x)=\tanh\left(\frac{e\cdot x-b}{\sqrt2}\right),
 \qquad |e|=1.
\]
In particular this conclusion holds for stable solutions on $\mathbb R^3$ with $E_1(u;B_R)\leq CR^2$ for $R\geq1$.
\end{theorem}
\begin{proof}
A constant solution is one of the two wells or zero, and zero is unstable as in Lemma \ref{lem:phase-stable-ode}. For a nonconstant solution, Proposition \ref{prop:phase-quadratic-flatness} and Lemma \ref{lem:phase-flat-to-onedimensional} reduce it to one dimension. Stability passes to the profile, and Lemma \ref{lem:phase-stable-ode} classifies it. The energy bound in the three-dimensional case implies $\int_{B_R}|Du|^2\leq2CR^2$. This proves the version of \cite{ChAC}, Theorem 5.10 and Section 5.4, with constant wells included and the unique-continuation dependency explicit.
\end{proof}

\begin{corollary}[Classification of bounded stable solutions in the plane]\label{cor:phase-planar-stable}
Every stable solution $u:\mathbb R^2\to[-1,1]$ of $\Delta u=W'(u)$ is a constant well or a planar heteroclinic. Therefore every bounded strictly monotone entire solution in the plane is a heteroclinic.
\end{corollary}
\begin{proof}
Lemma \ref{lem:phase-entire-estimates} gives a uniform gradient bound, so $\int_{B_R}|Du|^2\leq CR^2$. Apply Theorem \ref{thm:phase-stable-growth-rigidity}. Strict monotonicity gives stability by Corollary \ref{cor:phase-monotone-stability} and excludes constant solutions. This is the two-dimensional classification of \cite{ChAC}, Theorem 5.9 and Section 5.3, using the stated analytic dependencies.
\end{proof}
